The descriptive pattern is consistent with earlier observations that
averaging can provide a stronger relative benefit while a learning-rate
policy maintains stochastic updates, and a smaller difference once the raw
model has caught up during annealing \cite{MoralesBrotons2024,Ajroldi2025}.
The newer benchmark already studies this interaction directly. The present
results add a transparent small-model diffusion example with retained paired
states; they do not establish a new general principle.

The experiment intervenes on the full learning-rate path. Changes in the
terminal rate, the sequence of updates and the parameters entering the EMA
occur together. A causal explanation restricted to the last rate would
require another design. Likewise, the four-seed duration and interaction
intervals do not support a reliable conclusion about how the schedule
contrast changes with longer training.

Only two fixed EMA decays were tested. Architecture-dependent and
training-dependent averaging profiles \cite{Karras2024}, uniform averaging
\cite{Izmailov2018}, and switching averaged parameters back into training
\cite{Li2024} are different choices. The retained data cannot rank them or
justify a universal decay recommendation. The contrary ordinary-EMA
outcomes in the literature also argue against treating an improvement as a
condition for a valid research package.
